\section{Vision}
\label{sec:vision}

When it comes to data (and for that matter, any kind of data, not just multi-model), what we really want is to be able to extract information from it---i.e., ask questions about it and get answers---in a process commonly known as \term{querying}. In the context of multi-model data, this process naturally extends across boundaries of different models and, therefore, database systems.

Besides this primary goal, there are several other milestones that will help us to achieve it. We need to be able to \term{model} data so that we can reason about its structure and semantics---a necessary prerequisite for querying and some of the following milestones. Sometimes, the data is modeled first and created later. In other cases, we want to \term{infer} the schema from the data.

Both modeling and inference are closely related to the process of \term{evolution}. Because if the schema of the data changes (and in modern applications it is usually the case of \uv{when} instead of \uv{if}), we have to propagate the changes both to the data and to the queries that depend on it. In a multi-model data world, this might require us to \term{transform} the data between different models.

Lastly, our motivation is not only about an \emph{ability} to do something, but also about the \emph{efficiency} of doing so. In accordance with our primary goal, we want to be able to query data as efficiently as possible. In this regard, multi-model data offers a unique opportunity to optimize querying by choosing the most appropriate model for the task at hand. This is very difficult to do manually, so we want the system to be able to adapt itself---in a process called \term{self-adaptation}---to the current workload and data characteristics.

\subsection{Level of Automation}
\label{sec:vision-automation}

All of the above-described capabilities can be achieved with a different amount of automation. More automation is generally better, as it reduces the amount of manual work and expertise required of users. It also allows the system to find solutions that might not be obvious to users, or that would be too time-consuming to do manually. Most importantly, an automated process is reliable and does not make errors.

The price for this is increased complexity of the system and potentially reduced control over the process. Therefore, we automate the capabilities gradually following three distinct levels~\citep{ideas2022vision}:
\begin{itemize}
    \item \term{User-specified changes} -- The users change the schema of the data. E.g., they add a column \object{deletedAt} to a table \object{review}, allowing them to track when reviews are deleted. Or, the users decide to transform the data from the columnar model to the graph model, as the latter is more suitable for the current workload.
    \item \term{Data-driven changes} -- The system detects that the data has changed its structure, e.g., by observing that a new column \object{deletedAt} has been added to the \object{review} table, and updates the schema. Or, it finds out that labels (matching existing tables in the columnar database) are now present in the graph database, and creates the corresponding logical schemas.
    \item \term{Self-adaptation} -- The system proactively analyzes the workload and data characteristics, and decides to transform the data from a columnar model to a graph model (similarly how a database expert would do), without any user intervention.
\end{itemize}
These points focus only on the \emph{source of a change}, not on the \emph{way it is propagated}. The latter is covered in~\Cref{sec:evolution-automation}.

\subsection{Research Goals}

The previous paragraphs provide a high-level overview of the motivation for this thesis. For a more in-depth description, see~\paperRef{vldb_2024}. With this in mind, we can now formulate the concrete research goals of our overall approach. The first three goals were already addressed by previous works~\citep{mm_cat, mm_cat_journal, mm_quecat, koupil2023universal}:
\begin{itemize}
    \item \term{Modeling} -- How can heterogeneous data models be represented in a common semantic framework without erasing their model-specific structure?
    \item \term{Transformations} -- How can data be transformed between different models without the need to write custom transformation rules for each pair of models?
    \item \term{Querying} -- How can queries be unified across different data models, supporting all features of each model, and allowing for cross-model queries and optimizations?
\end{itemize}
The rest of the goals are novel and thus are the main focus of this thesis:
\begin{itemize}
    \item \term{Evolution} -- How can changes be propagated across conceptual and logical schemas, data, and queries in a multi-model environment? How can uncertain or ambiguous evolution decisions be supported by data-derived semantic evidence?
    \item \term{Self-adaptation} -- How can a multi-model system choose suitable representations for a workload without exhaustively materializing and benchmarking all alternatives?
    \item \term{Benchmarking} -- How can such systems be evaluated when real multi-model benchmarks and workloads are scarce?
    \item \term{Limits of the framework} -- What are the limits of the proposed framework, and which parts require alternative approaches?
\end{itemize}

State-of-the-art approaches to data management provide only partial solutions to these challenges. Just to name a few examples of their limitations, they typically:
\begin{itemize}
    \item support multiple data models, but without a deeper unified semantic foundation,
    \item allow transformations and evolution, but within a single data model,
    \item optimize physical design, but not decisions like \uv{relational vs. document vs. graph database}, and
    \item support querying, but not long-term evolution (and vice versa).
\end{itemize}

As we can see, what is missing is a principled management layer that connects modeling, querying, evolution, and self-adaptation. Developing such a layer is the main goal of our research, and its extension to the above-described aspects of multi-model data management is the main contribution of this thesis.

\todo{zmínit i "necíle"?}

\subsection{Challenges}

\todo{Možná by nebylo od věci někde na 1-2 stránky shrnout ty základní modely, v čem se liší a tudíž proč je těžké řešit je všechny najednou.}

\todo{jedna DB nemusí stačit (i když je MM), můžeme chtít jiné škálování

!zmínit redundancy
}
